\section{Dataset}
\label{sec:dataset}

VIGIL-Dataset is a curated benchmark of AI-related copyright and intellectual-property litigation designed for structured risk analysis.

\subsection{Construction and Statistics}
\label{sec:construction}

The dataset compiles CourtListener/RECAP filings and supplementary legal-news signals spanning January 2020 to December 2025, with the majority of filings concentrated in 2023--2025.
Cases were obtained via a retrospective keyword search of RECAP complaints and were manually screened to confirm that AI development, training, or deployment was at issue.
The scope is intentionally broader than copyright-only trackers: it includes privacy, publicity, product-liability, and contract theories when AI systems are material to the dispute, and counts each docket separately (including related and consolidated filings).
One case corresponds to one docket; consolidated actions are therefore counted per docket.
Table~\ref{tab:dataset} summarizes the final annotated subset (1{,}247 dockets after screening and dedup; $\approx$980 distinct disputes).
Public AI-copyright trackers typically list a few hundred distinct disputes.
Our 1{,}247 figure counts dockets, not disputes: consolidated and related actions are retained as separate records.
After merging on dispute identity the unique-dispute count is approximately 980.
The top-three defendants account for 41.3\% of dockets, reflecting the concentration of high-profile generative-AI litigation; we release the defendant-frequency table so that users can assess applicability to smaller vendors.

\begin{table}[t]
\centering
\small
\begin{adjustbox}{max width=\columnwidth}
\begin{tabular}{lc}
\toprule
Attribute & Value \\
\midrule
Total cases & 1{,}247 \\
Total documents & 8{,}934 \\
Coverage period & 2020--2025 \\
Unique defendants & 43 \\
Risk labels & 3 (High/Med/Low) \\
Status: Initial / Active & 298 / 695 \\
Status: Settled / Dismissed & 148 / 106 \\
Product-mapping GT cases & 312 \\
\bottomrule
\end{tabular}
\end{adjustbox}
\caption{VIGIL-Dataset summary statistics.
Status counts sum to 1{,}247.}
\label{tab:dataset}
\end{table}

\subsection{Annotation and Availability}
\label{sec:annotation}

An automated pre-labeler first assigns a triage score using an additive scoring matrix (supplementary material).
Three legal research specialists then independently review a stratified sample of 400 cases under a fixed guideline.
During adjudication, annotators see docket text, product attributes, candidate theory cards, and query dates, but neither system identity nor the pre-labeler's score.
The remaining 847 cases retain automated pre-labels; a 5\% spot-check yields 91\% agreement with at least one specialist.
Inter-annotator agreement on the 400-case sample is mean pairwise Cohen's $\kappa = 0.74$ (Fleiss' $\kappa = 0.72$).
Because retained labels can inflate metrics, we release per-record provenance, restrict all primary claims to the human-only evaluation, and treat retained-label results solely as scale-oriented comparisons.
The dataset, guidelines, split files, seeds, and evaluation scripts are supplied as anonymized supplementary material.
